\documentclass[12pt, a4paper]{article}

% -----------------------------------------------------------
% Packages
% -----------------------------------------------------------

\usepackage[utf8]{inputenc}    % UTF-8 encoding
\usepackage[margin=1in]{geometry}          % Page geometry
\geometry{margin=1in}          % 1-inch margins
\usepackage{amsmath, amssymb} % Math symbols and equations
\usepackage{graphicx}          % Include images
\usepackage{hyperref}          % Clickable links in PDF
\usepackage{setspace}          % Line spacing
\usepackage{xcolor}            % Text color
\usepackage{titlesec}          % Custom section titles
\usepackage{enumitem}          % Custom lists

% -----------------------------------------------------------
% Document Configuration
% -----------------------------------------------------------

\title{Fisica}
\author{Claudio Tessa}
\date{\the\year}
    
% Custom section styling
\titleformat{\section}
  {\normalfont\large\bfseries}{\thesection}{1em}{}

% -----------------------------------------------------------
% Begin Document
% -----------------------------------------------------------

\begin{document}
\maketitle

\section*{Informazioni generali}
L'esame verrà valutato con una prova scritta ed una prova orale facoltativa. La prova scritta consisterà in alcuni esercizi da svolgere e domande di carattere teorico sugli argomenti trattati. La valutazione 30 e lode è raggiungibile solo con la prova orale (il massimo voto accessibile con la sola prova scritta è 30).

\tableofcontents
\newpage

% -----------------------------------------------------------
% Introduzione
% -----------------------------------------------------------

\section{Introduzione}

La fisica studia i fenomeni naturali ricercando le leggi che li governano. La fisica "antica" era intrinsecamente legata alla filosofia e religione, e vigeva il principio di autorità. La fisica moderna nasce con Galileo: abbiamo la separazione della realtà in ciò che è oggettivo, quindi misurabile misurabile, e ciò che non lo è, che quindi è ritenuto soggettivo.\\
I misurabili della realtà possono essere descritti mediante la matematica, e permettono la realizzazione di una (i) ipotesi sul fenomeno, di una (ii) teoria che spieghi il fenomeno stesso e di (iii) dedurre le conseguenze.\\
Chiunque, conoscendo le condizioni di esecuzione può, avendone i mezzi e le conoscenze necessarie, essere capace di ripetere l'esperimento

% -----------------------------------------------------------
% Cinematica
% -----------------------------------------------------------

\section{Cinematica}
La cinematica è parte della Meccanica e si occupa della descrizione spazio-temporale del moto dei corpi, assumendo che esso avvenga con continuità, indipendentemente dalle cause del moto. \\
Ogni moto esiste solo all'interno di un sistema di riferimento, la cui scelta lo fa apparire semplice o complesso. Poichè un corpo non cambia a seconda del sistema di riferimento, la scelta di tale sistema è fondamentale per semplificare al massimo la comprensione e formulazione del fenomeno stesso

\subsection{Cinematica del punto materiale (1D)}
Sia dato un sistema di riferimento orientato $x$, e un punto materiale (un oggetto puntiforme di massa $m$). Abbiamo che:
\begin{itemize}
    \item \textbf{Spostamento}: $s = \Delta x = x_f - x_i$
    \item \textbf{Velocità media}: $v_m =\frac{x_f - x_i}{t_f - t_i} = \frac{\Delta x}{\Delta t}$
    \item \textbf{Velocità istantanea}: $v_x = \lim_{\Delta t \to 0}{\frac{\Delta x}{\Delta t}} = \frac{dx}{dt}$
    \item \textbf{Accelerazione media}: $a_m = \frac{v_f - v_i}{t_f - t_i}= \frac{\Delta v}{\Delta t}$
    \item \textbf{Accelerazione istantanea}: $a_x = \lim_{\Delta t \to 0}{\frac{\Delta v}{\Delta t}} = \frac{d^2x}{d^2t}$ 
\end{itemize}

\noindent
Se l'accelerazione varia continuamente, il moto non è facile da analizzare. Lo studio dei moti si semplifica quando il moto ha alcune proprietà:
\begin{itemize}
    \item \textbf{Moto rettilineo uniforme}: velocità costante e accelerazione zero. \\
        $x(t) = x_i + v_x \cdot t$\\
        $v_x = k$, con $k \in \mathbb{R}$
    \item \textbf{Moto con accelerazione costante (o uniformemente accelerato)}: accelerazione istantanea uguale all'accelerazione media in ogni istante. La velocità cambia linearmente durante tutto il moto.\\
        $x(t) = x_i + v_x \cdot t + \frac{1}{2} a_x \cdot t^2$ \\
        $v_x (t) = v_{ix} + a_x \cdot t$ \\
        $a_x = k$, con $k \in \mathbb R$
\end{itemize}


% -----------------------------------------------------------
% Figures
% -----------------------------------------------------------

\section{Adding Figures}
You can include images using the \texttt{graphicx} package:
\begin{figure}[h!]
	\centering
	\includegraphics[width=0.5\textwidth]{example-image} % Replace with your image
	\caption{An example figure.}
	\label{fig:example}
\end{figure}

% -----------------------------------------------------------
% Lists
% -----------------------------------------------------------

\section{Lists}
Here is an example of an itemized list:
\begin{itemize}
	\item First item
	\item Second item
	\item Third item
\end{itemize}

And an enumerated list:
\begin{enumerate}
	\item Step one
	\item Step two
	\item Step three
\end{enumerate}

% -----------------------------------------------------------
% Conclusion
% -----------------------------------------------------------

\section{Conclusion}
This is a simple document to get started with LaTeX. Explore adding more sections, tables, references, and advanced formatting!

% -----------------------------------------------------------
% End Document
% -----------------------------------------------------------
\end{document}

